% resumen.tex
\chapter*{Resumen}
\addcontentsline{toc}{chapter}{Resumen}
\begin{abstract}
El diseño de placas base para columnas de acero es uno de los elementos más
importantes en la seguridad estructural de edificaciones, ya que estas
conexiones transmiten las cargas axiales, cortantes y momentos flexionantes de
la superestructura hacia la cimentación de concreto. Las Normas Técnicas
Complementarias para el Diseño y Construcción de Estructuras de Acero de la
Ciudad de México (NTC-DCEA) fueron actualizadas en 2023, incorporando los
factores de confinamiento por estribos $R_{et}$ y $R_{ev}$ para los estados
límite de resistencia al cono en tensión y cortante, así como criterios más
restrictivos de distancia mínima al borde libre para anclas de alta
resistencia. El sismo del 19 de septiembre de 2017 evidenció fallas frágiles
por arrancamiento de cono de concreto en conexiones columna-cimentación
diseñadas conforme a las NTC-DCEA 2017, lo que motivó estas actualizaciones.

La presente investigación desarrolla un análisis paramétrico numérico
comparativo entre ambas normativas, ejecutando tres barridos sistemáticos que
cubren las variables de mayor divergencia: la profundidad de empotramiento
$h_{ef}$ y el espaciado de estribos del dado, el diámetro de ancla y su
influencia en los requisitos geométricos del pedestal, y la resistencia del
concreto $f'_c$ en combinación con el espaciado del refuerzo transversal. Los
resultados, sintetizados en cuatro figuras de análisis, demuestran que la
NTC-DCEA 2023 es consistentemente más permisiva que la NTC-DCEA 2017 en la
resistencia al cono ---al incorporar $R_{et}$ y $R_{ev}$--- y más restrictiva
en los requisitos geométricos para anclas A325 y A449, con implicaciones
directas en el dimensionamiento de pedestales y en la instrumentación
práctica del Diseño por Capacidad en estructuras sismorresistentes.

\textbf{Palabras clave:} placas base, conexiones columna-cimentación,
NTC-DCEA, diseño de anclajes, análisis paramétrico normativo.
\end{abstract}
